% 3-7-1-lightgbm.tex
%  Budget: 3pp.  Converted from thesis/drafts/3-7-1-lightgbm.md.
%  Source map: gradient boosting on trees as a standard non-linear baseline in
%   recent EPF work -> jiang_review_2018 , oconnor_review_2025 ;
%   the per-hour decomposition mirrors the LEAR reference -> lago_forecasting_2021 .
%  Hyperparameters verified 2026-09-05 against configs/tuned/lightgbm_params.yaml:
%   n_estimators 963, lr 0.024499, num_leaves 15, min_child_samples 11,
%   feature_fraction 0.542, bagging_fraction 0.513, bagging_freq 3,
%   lambda_l1 0.31636, lambda_l2 7.04e-08, validation_mae 3.7316.
%  Test-period numbers from results_canonical.csv: MAE 3.968, rMAE 0.435.
%  OOD number from ood_stress.csv: rMAE 1.828.
%  LightGBM is already footnoted in ch1-2 -> \lr{} only, no footnote.

\subsection{\lr{LightGBM}}
\label{sec:3-7-1}

نخستین مدل غیرخطی این پژوهش \lr{LightGBM} است: پیاده‌سازی کارآمدی از تقویت گرادیانی
بر درخت‌های تصمیم، که در ادبیات پیش‌بینی قیمت برق به‌عنوان نماینده‌ی استاندارد
خانواده‌ی مدل‌های درختی به کار می‌رود \cite{jiang_review_2018, oconnor_review_2025}.
انتخاب آن \lr{—} به‌جای جنگل تصادفی یا \lr{XGBoost} که هر دو آگاهانه از دامنه‌ی پژوهش
کنار گذاشته شدند \lr{—} بر دو ملاحظه استوار است: کارایی محاسباتی، که واسنجی روزانه
را در ۷۲۸ مبدأ ممکن می‌سازد، و توان بالای آن در داده‌های جدولی با ویژگی‌های ناهمگن،
که دقیقاً ساختار ماتریس ویژگیِ \cref{sec:3-4} است.

\subsubsection*{ساختار مدل}

به‌ازای هر ساعت شبانه‌روز یک رگرسور مستقل برازش می‌شود؛ یعنی ۲۴ مدل مجزا که هیچ‌یک
از دیگری اطلاعی ندارد. این تصمیم عامدانه با ساختار ساعت‌به‌ساعتِ \lr{LEAR-LASSO} و
\lr{SARIMAX} (\cref{sec:3-6}) هم‌قرینه است، تا اختلاف عملکرد میان مدل خطی و مدل درختی به
خودِ خانواده‌ی مدل بازگردد و نه به تفاوت در نحوه‌ی تجزیه‌ی افق پیش‌بینی.

هر یک از این ۲۴ مدل، کل ماتریس ویژگیِ مشترک را دریافت می‌کند: هر ۲۴۷ ستون، شامل
وقفه‌های قیمتی، وقفه‌ها و مقادیر روزپیشِ متغیرهای برون‌زا، و متغیرهای مجازیِ روز
هفته. هیچ مسیر ویژگی‌سازیِ جداگانه‌ای برای این مدل وجود ندارد و در نتیجه سطح دومی
برای نشت اطلاعات ایجاد نمی‌شود.

یادداشت پیاده‌سازی: پیمانه‌ی این مدل عامدانه نامی متفاوت از خودِ بسته‌ی \lr{LightGBM}
دارد، تا در هیچ شرایطی با آن اشتباه گرفته نشود یا آن را سایه نیندازد.

\subsubsection*{آهنگ واسنجی}

برخلاف \lr{SARIMAX} که برازش کامل آن در هر مبدأ عملی نبود (\cref{sec:3-6})، \lr{LightGBM}
بر پنجره‌ی واسنجیِ ۱۰۹۲روزه در چند ثانیه برازش می‌شود. از این رو هیچ مصالحه‌ای در آهنگ
واسنجی لازم نیست و مدل در هر مبدأ به‌طور کامل بازبرازش می‌گردد \lr{—} یعنی دقیقاً
همان پروتکل واسنجیِ روزانه‌ای که مدل مرجع \lr{LEAR} از آن پیروی می‌کند
\cite{lago_forecasting_2021}. این هم‌ترازی، مقایسه‌ی مستقیم دو مدل را از شبهه‌ی «تفاوت
در تازگی برازش» پاک می‌کند.

\subsubsection*{تنظیم ابرپارامترها}

ابرپارامترها با ۵۰ آزمایش \lr{Optuna} تنظیم شده‌اند. پنجره‌ی اعتبارسنجی، از ۵ ژانویه‌ی
۲۰۱۵ تا ۳ ژانویه‌ی ۲۰۱۶، اکیداً پیش از دوره‌ی آزمون قرار دارد، و پنجره‌ی برازشِ آن، از ۱۶
ژانویه‌ی ۲۰۱۲ تا ۴ ژانویه‌ی ۲۰۱۵، پیش از خودِ پنجره‌ی اعتبارسنجی است. بدین ترتیب هیچ
مشاهده‌ای از دوره‌ی آزمون در انتخاب ابرپارامترها نقشی ندارد. بذر نمونه‌بردارِ جست‌وجو
نیز ۴۲ تثبیت شده تا خودِ فرایند تنظیم بازتولیدپذیر بماند.

تابع هدف بر زیان قدرمطلق تنظیم شده است تا زیانِ آموزش با معیار سرشناسِ گزارش
هم‌راستا باشد؛ همین انتخاب برای \lr{LSTM} نیز تکرار می‌شود تا مقایسه‌ی آن دو مخدوش
نشود. بهترین پیکربندیِ یافت‌شده، با \lr{MAE} اعتبارسنجیِ \lr{3.7316}، چنین است:
شمار درخت‌ها ۹۶۳، نرخ یادگیری \lr{0.0245}، شمار برگ‌ها ۱۵، کمینه‌ی نمونه در
برگ ۱۱، نسبت ویژگی \lr{0.542}، نسبت نمونه‌گیری \lr{0.513} با تناوب ۳، و
جریمه‌های \lr{L1} و \lr{L2} به‌ترتیب \lr{0.316} و \lr{7.04e-08}.

جهت‌گیریِ این جست‌وجو خود گویاست. نسبت به مقادیر پیش‌فرض، شمار برگ‌ها از ۶۳ به ۱۵
کاهش و در مقابل شمار درخت‌ها از ۵۰۰ به ۹۶۳ افزایش یافته و نرخ یادگیری نصف شده
است. یعنی جست‌وجو به سمت مدلی کوچک‌تر، منظم‌تر و آهسته‌تر حرکت کرده \lr{—} الگویی که
با ماهیت پرنوفه‌ی سری قیمت برق و خطر بیش‌برازش بر ویژگی‌های فراوان سازگار است.

\subsubsection*{تضمین بازتولیدپذیری}

سه پارامتر، فارغ از محتوای فایل پیکربندی یا نتیجه‌ی تنظیم، در خودِ کد تحمیل می‌شوند:
بذر تصادفی ۴۲، پرچم جبرگرایی، و مهم‌تر از هر دو، شمار ثابت نخ‌های پردازش.

علت تثبیت صریح شمار نخ‌ها \lr{—} به‌جای واگذاری آن به مقدار خودکار \lr{—} ظریف و
ضروری است: تضمین جبرگراییِ این کتابخانه تنها به‌ازای شمار نخِ یکسان برقرار است. اگر
این مقدار با مشخصات ماشین شناور بماند، اجرای همان کد بر سخت‌افزاری دیگر می‌تواند
اعدادی متفاوت تولید کند، و آنگاه ادعای بازتولیدپذیریِ نتایج تثبیت‌شده بی‌اعتبار
می‌شود. این مقدار پس از آنکه آزمایش‌های تک‌نخیِ تنظیم به درازا کشید افزایش یافت، و
تغییر آن به‌عنوان یک انتخابِ ثبت‌شده و نه یک تنظیم محیطی، در دفتر تصمیمات درج شده
است.

\subsubsection*{عملکرد و ارجاع به فصل چهارم}

بر دوره‌ی آزمونِ ۷۲۸ مبدأیی، \lr{MAE} ساعتی این مدل \lr{3.968} و \lr{rMAE} آن
\lr{0.435} است؛ جدول کامل در \cref{sec:4-2} می‌آید. دو ملاحظه‌ی صادقانه باید همین‌جا ثبت شود.

نخست آنکه این مدل غیرخطی، مدل خطیِ \lr{LEAR-LASSO} را در \lr{MAE} ساعتی پشت سر
نمی‌گذارد: \lr{3.968} در برابر \lr{3.899}. برتریِ روش‌های انعطاف‌پذیر بر مدل خطیِ
مرجع، در این بازار و این دوره محقق نشده است.

دوم آنکه در آزمون برون‌توزیعیِ فصل پنجم، همین مدل بیشترین افت را در میان تمام مدل‌ها
نشان می‌دهد، با \lr{rMAE} برابر \lr{1.828}. تفسیر این دو یافته در کنار یکدیگر
\lr{—} یعنی اینکه انعطاف‌پذیریِ بالاتر در درون‌دوره مزیتی نساخت اما در برون‌توزیع
هزینه‌ی سنگین‌تری تراشید \lr{—} در \cref{sec:5-2} می‌آید.
